\markdownRendererDocumentBegin
\markdownRendererWarning{The `hybrid` option has been soft-deprecated.}{The `hybrid` option has been soft-deprecated.}{Consider using one of the following better options for mixing TeX and markdown: `contentBlocks`, `rawAttribute`, `texComments`, `texMathDollars`, `texMathSingleBackslash`, and `texMathDoubleBackslash`. For more information, see the user manual at <https://witiko.github.io/markdown/>.}{Consider using one of the following better options for mixing TeX and markdown: `contentBlocks`, `rawAttribute`, `texComments`, `texMathDollars`, `texMathSingleBackslash`, and `texMathDoubleBackslash`. For more information, see the user manual at <https://witiko.github.io/markdown/>.}\markdownRendererSectionBegin
\markdownRendererHeadingOne{Teoría de Números}\markdownRendererInterblockSeparator
{}\markdownRendererSectionBegin
\markdownRendererHeadingTwo{Números Primos y Factores}\markdownRendererInterblockSeparator
{}\markdownRendererStrongEmphasis{Definición:}\markdownRendererHardLineBreak
{}Un número $a$ es \markdownRendererStrongEmphasis{factor} (o \markdownRendererStrongEmphasis{divisor}) de $b$ si $a$ divide a $b$, denotado $a \mid b$.\markdownRendererHardLineBreak
{}Un entero $n > 1$ es \markdownRendererStrongEmphasis{primo} si sus únicos divisores positivos son $1$ y $n$.\markdownRendererParagraphSeparator
{}\markdownRendererStrongEmphasis{Factorización prima única:}\markdownRendererHardLineBreak
{}Todo entero $n>1$ se puede escribir de forma única (hasta orden) como:\markdownRendererSoftLineBreak
{}$$\markdownRendererSoftLineBreak
{}n = p_1^{\alpha_1} p_2^{\alpha_2} \cdots p_k^{\alpha_k}\markdownRendererSoftLineBreak
{}$$\markdownRendererParagraphSeparator
{}$\clearpage$\markdownRendererInterblockSeparator
{}\markdownRendererSectionBegin
\markdownRendererHeadingThree{Cantidad de factores}\markdownRendererInterblockSeparator
{}Sea $\tau(n)$ el número de factores de un entero n. Por ejemplo, $\tau(12) = 6$, porque los factores de 12 son 1, 2, 3, 4, 6 y 12. Para calcular el valor de $\tau(n)$, podemos usar la fórmula\markdownRendererParagraphSeparator
{}$$\markdownRendererSoftLineBreak
{}\tau(n) = \prod_{i=1}^{k} (\alpha_i + 1)\markdownRendererSoftLineBreak
{}$$\markdownRendererParagraphSeparator
{}ya que para cada primo $p_i$, hay $\alpha_i + 1)$ formas de elegir cuántas veces aparece en el factor. Por ejemplo, como $12 = 2^2 \cdot 3^1$ $\tau(12) = (2+1) \cdot 1(2) = 6$.\markdownRendererParagraphSeparator
{}\markdownRendererStrongEmphasis{Complejidad:} $O(\sqrt{n})$\markdownRendererInterblockSeparator
{}\markdownRendererInputFencedCode{./_markdown_main/f77317809c9f28d2c621be39820ed13e.verbatim}{cpp}{cpp}\markdownRendererInterblockSeparator
{}$\clearpage$\markdownRendererInterblockSeparator
{}
\markdownRendererSectionEnd 
\markdownRendererSectionEnd \markdownRendererSectionBegin
\markdownRendererHeadingTwo{Suma de factores}\markdownRendererInterblockSeparator
{}Sea $\sigma(n)$ la suma de los factores de un entero n. Por ejemplo, $\sigma(12) = 28$, porque (1 + 2 + 3 + 4 + 6 + 12 = 28). Para calcular el valor de $\sigma(n)$, podemos usar la fórmula\markdownRendererParagraphSeparator
{}$$\markdownRendererSoftLineBreak
{}\sigma(n) = \prod_{i=1}^{k} (1 + p_i + \dots + p_i^{\alpha_i}) = \prod_{i=1}^{k} \frac{p_i^{\alpha_i+1} - 1}{p_i - 1}\markdownRendererSoftLineBreak
{}$$\markdownRendererParagraphSeparator
{}donde la última forma se basa en la fórmula de progresión geométrica.\markdownRendererSoftLineBreak
{}Por ejemplo, $ \sigma(12) = \frac{2^3 - 1}{2 - 1} \cdot \frac{3^2 - 1}{3 - 1} = 28$\markdownRendererParagraphSeparator
{}\markdownRendererStrongEmphasis{Complejidad:} $O(\sqrt{n})$\markdownRendererInterblockSeparator
{}\markdownRendererInputFencedCode{./_markdown_main/8655855bb1415a3150e01e31d5421ee9.verbatim}{cpp}{cpp}\markdownRendererInterblockSeparator
{}$\clearpage$\markdownRendererInterblockSeparator
{}\markdownRendererSectionBegin
\markdownRendererHeadingThree{Lista de factores}\markdownRendererInterblockSeparator
{}Observa que cada factor primo aparece en el vector tantas veces como divide el número.\markdownRendererParagraphSeparator
{}Por ejemplo $12 = 2^2 * 3$, por lo que el resultado de la función es $[2,2,3]$\markdownRendererParagraphSeparator
{}\markdownRendererStrongEmphasis{Complejidad:} $O(\sqrt{n})$\markdownRendererInterblockSeparator
{}\markdownRendererInputFencedCode{./_markdown_main/51cbb90f5e3949a15bb8a934fea7eba5.verbatim}{cpp}{cpp}\markdownRendererInterblockSeparator
{}Variante que toma la criba y devuelve un map donde el valor de cada clave es la potencia, ejemplo\markdownRendererSoftLineBreak
{}$12 = {2:2, 3:1}$\markdownRendererParagraphSeparator
{}\markdownRendererStrongEmphasis{Complejidad:} $O(\pi(\sqrt{n}) + \log n) \approx O(\frac{\sqrt{n}}{\log n} + \log n)$\markdownRendererInterblockSeparator
{}\markdownRendererInputFencedCode{./_markdown_main/194cb71f0d134dfc899d72cd6af56d71.verbatim}{cpp}{cpp}\markdownRendererInterblockSeparator
{}
\markdownRendererSectionEnd 
\markdownRendererSectionEnd \markdownRendererSectionBegin
\markdownRendererHeadingTwo{Pruebas de primalidad}\markdownRendererInterblockSeparator
{}Revisar si un número es primo.\markdownRendererInterblockSeparator
{}\markdownRendererInputFencedCode{./_markdown_main/f4f0bf5f532455a33d48a2536233211b.verbatim}{cpp}{cpp}\markdownRendererInterblockSeparator
{}\markdownRendererSectionBegin
\markdownRendererHeadingThree{Algoritmos probabilisticos}\markdownRendererInterblockSeparator
{}Estos algorimtos se basan en comprobar de manera probabilistica basado en el Teoerema debil de Fermat, cuanto mayor sea la cantidad de iteraciones, mneor es la probabilidad de equviocarse.\markdownRendererInterblockSeparator
{}\markdownRendererSectionBegin
\markdownRendererHeadingFour{Prueba de Fermat}\markdownRendererInterblockSeparator
{}La prueba de Fermat puede fallar con los números de Charmichael  ($a^{n-1} \equiv 1 \bmod n$ se cumple para todo$a$ coprimo con $n$, por ejemplo $561 = 3 \times 11 \times 17$), sin embargo en la practica se sigue usando ya que es bastante rápida y los números de Charmichael son raros.\markdownRendererInterblockSeparator
{}\markdownRendererInputFencedCode{./_markdown_main/9423d86fae1e1ffd6badcb31f712cab6.verbatim}{cpp}{cpp}\markdownRendererInterblockSeparator
{}
\markdownRendererSectionEnd \markdownRendererSectionBegin
\markdownRendererHeadingFour{Prueba de Miller}\markdownRendererInterblockSeparator
{}Miller desmostro que el algoritmo se puede volver determinista si se prueban todas bases $a \leq 2 ln(n)^2$,\markdownRendererSoftLineBreak
{}siendo que para los números de 64 bit, basta con revisar las primeras 12 bases primas.\markdownRendererParagraphSeparator
{}\markdownRendererInputFencedCode{./_markdown_main/70f32b8a08478145d8b534c93c77e01e.verbatim}{cpp}{cpp}\markdownRendererInterblockSeparator
{}$\clearpage$\markdownRendererInterblockSeparator
{}
\markdownRendererSectionEnd 
\markdownRendererSectionEnd \markdownRendererSectionBegin
\markdownRendererHeadingThree{Primos como diferencia de cuadrados}\markdownRendererParagraphSeparator
{}Un numero impar se puede escribir $n = pq$ como la diferencia de cuadrados $n = a^2 - b^2$.\markdownRendererParagraphSeparator
{}$$n = (\frac{p+q}{2})^2 - (\frac{p-q}{2})^2$$\markdownRendererParagraphSeparator
{}Estos valores $p$ y $q$ se pueden obtener con el metodo de factorización de Fermat. $O(|p-q|)$ cabe\markdownRendererSoftLineBreak
{}mencionar que si los números son demasiado grande o la distancia entre $p$ y $q$, este algoritmo es\markdownRendererSoftLineBreak
{}extremadamente lento.\markdownRendererInterblockSeparator
{}\markdownRendererInputFencedCode{./_markdown_main/a9a4b39bbf10fd16685c491bd77951e3.verbatim}{cpp}{cpp}\markdownRendererInterblockSeparator
{}
\markdownRendererSectionEnd 
\markdownRendererSectionEnd \markdownRendererSectionBegin
\markdownRendererHeadingTwo{Computar números primos}\markdownRendererParagraphSeparator
{}\markdownRendererSectionBegin
\markdownRendererHeadingThree{Criba de Erastotenes}\markdownRendererInterblockSeparator
{}\markdownRendererStrongEmphasis{Idea:} marcar múltiplos de cada primo hasta $n$.\markdownRendererHardLineBreak
{}\markdownRendererStrongEmphasis{Complejidad:} $O(n \log \log n)$\markdownRendererInterblockSeparator
{}\markdownRendererInputFencedCode{./_markdown_main/62198306d6a4f3c98cb7c8f3275e0616.verbatim}{cpp}{cpp}\markdownRendererInterblockSeparator
{}$\clearpage$\markdownRendererInterblockSeparator
{}
\markdownRendererSectionEnd \markdownRendererSectionBegin
\markdownRendererHeadingThree{Criba Lineal}\markdownRendererInterblockSeparator
{}\markdownRendererStrongEmphasis{Idea:} Existe una optimización que permite realizar la criba en $O(n)$ a cambio de usar más memoria, que solo debe ser usada si $N \leq 10^7$\markdownRendererSoftLineBreak
{}\markdownRendererStrongEmphasis{Complejidad} : $O(n)$\markdownRendererInterblockSeparator
{}\markdownRendererInputFencedCode{./_markdown_main/b9d19f4fd13f5a17a42e8a28995b2ec1.verbatim}{cpp}{cpp}\markdownRendererInterblockSeparator
{}
\markdownRendererSectionEnd \markdownRendererSectionBegin
\markdownRendererHeadingThree{Contar número de primos en un rango}\markdownRendererInterblockSeparator
{}\markdownRendererStrongEmphasis{Idea:} Realizar una suma prefijos\markdownRendererSoftLineBreak
{}\markdownRendererStrongEmphasis{Complejidad de precomputo:} $O(n \log \log \sqrt{n})$\markdownRendererSoftLineBreak
{}\markdownRendererStrongEmphasis{Complejidad de consulta:} $O(1)$\markdownRendererInterblockSeparator
{}\markdownRendererInputFencedCode{./_markdown_main/564d090d7d9420c98260c0d95094053f.verbatim}{cpp}{cpp}\markdownRendererInterblockSeparator
{}$\clearpage$\markdownRendererInterblockSeparator
{}
\markdownRendererSectionEnd \markdownRendererSectionBegin
\markdownRendererHeadingThree{Encontrar primos en un rango}\markdownRendererInterblockSeparator
{}Encontrar primos en un rango\markdownRendererParagraphSeparator
{}Encontrar números primos en un rango $[L, R]$ donde $R - L + 1 \approx 10^7$ y $R \leq 10^{12}$.\markdownRendererInterblockSeparator
{}\markdownRendererInputFencedCode{./_markdown_main/68cbc9644aa9ef02e1213c2627a4fd4e.verbatim}{cpp}{cpp}\markdownRendererInterblockSeparator
{}
\markdownRendererSectionEnd 
\markdownRendererSectionEnd \markdownRendererSectionBegin
\markdownRendererHeadingTwo{Algoritmo de Euclides}\markdownRendererParagraphSeparator
{}\markdownRendererSectionBegin
\markdownRendererHeadingThree{Máximo Comúm Divisior}\markdownRendererInterblockSeparator
{}$$\markdownRendererSoftLineBreak
{}\gcd(a,b) = \gcd(b, a \bmod b)\markdownRendererSoftLineBreak
{}$$\markdownRendererParagraphSeparator
{}\markdownRendererStrongEmphasis{Complejidad:} $O(\log \min(a,b))$.\markdownRendererInterblockSeparator
{}\markdownRendererInputFencedCode{./_markdown_main/8be859ccfb61681379fc273b4a7e9d4a.verbatim}{c}{c}\markdownRendererInterblockSeparator
{}
\markdownRendererSectionEnd \markdownRendererSectionBegin
\markdownRendererHeadingThree{Máximo Común Multiplo}\markdownRendererInterblockSeparator
{}$$\markdownRendererSoftLineBreak
{}\mathrm{lcm}(a,b) = \frac{a \cdot b}{\gcd(a,b)}\markdownRendererSoftLineBreak
{}$$\markdownRendererParagraphSeparator
{}\markdownRendererStrongEmphasis{Complejidad:} $O(\log \min(a,b))$.\markdownRendererInterblockSeparator
{}\markdownRendererInputFencedCode{./_markdown_main/394daf27a502d67ee061bed27f75b198.verbatim}{c}{c}\markdownRendererInterblockSeparator
{}$\clearpage$\markdownRendererInterblockSeparator
{}
\markdownRendererSectionEnd \markdownRendererSectionBegin
\markdownRendererHeadingThree{Euclides Extendido}\markdownRendererInterblockSeparator
{}\markdownRendererStrongEmphasis{Devuelve} $x,y$ tales que $a x + b y = \gcd(a,b)$ — útil para inversos modulares y diofánticas.\markdownRendererInterblockSeparator
{}\markdownRendererInputFencedCode{./_markdown_main/339a8dccaceea0fbfbb24fa4e78b65c8.verbatim}{cpp}{cpp}\markdownRendererInterblockSeparator
{}
\markdownRendererSectionEnd \markdownRendererSectionBegin
\markdownRendererHeadingThree{Numeros Coprimos}\markdownRendererInterblockSeparator
{}Dos enteros $a$ y $b$ son coprimos si $\gcd(a,b) = 1$.\markdownRendererInterblockSeparator
{}\markdownRendererInputFencedCode{./_markdown_main/62ab9e372f52b49dbcba35952dccf31d.verbatim}{cpp}{cpp}\markdownRendererInterblockSeparator
{}
\markdownRendererSectionEnd \markdownRendererSectionBegin
\markdownRendererHeadingThree{Funcion Totiente de Euler}\markdownRendererInterblockSeparator
{}La \markdownRendererStrongEmphasis{función totiente de Euler} $\varphi(n)$ da el número de enteros entre $1$ y $n$ que son coprimos con $n$ si\markdownRendererSoftLineBreak
{}$$\markdownRendererSoftLineBreak
{}n = p_1^{\alpha_1} \cdots p_k^{\alpha_k},\markdownRendererSoftLineBreak
{}$$\markdownRendererParagraphSeparator
{}$$\markdownRendererSoftLineBreak
{}\varphi(n) = \prod_{i=1}^{k} p_i^{\alpha_i - 1} (p_i - 1)\markdownRendererSoftLineBreak
{}$$\markdownRendererParagraphSeparator
{}\markdownRendererStrongEmphasis{Ejemplo:} $10 = 2 \cdot 5 ;\Rightarrow; \varphi(10) = 2^{1-1}(2-1) \cdot 5^{1-1}(5-1) = 4$\markdownRendererParagraphSeparator
{}\markdownRendererStrongEmphasis{Complejidad:} $O(\sqrt{n})$\markdownRendererInterblockSeparator
{}\markdownRendererInputFencedCode{./_markdown_main/0a84b81da84f0a4b46d3683d7dc800c4.verbatim}{cpp}{cpp}\markdownRendererInterblockSeparator
{}$\clearpage$\markdownRendererInterblockSeparator
{}
\markdownRendererSectionEnd \markdownRendererSectionBegin
\markdownRendererHeadingThree{Teorema de Euler}\markdownRendererInterblockSeparator
{}\markdownRendererStrongEmphasis{Teorema de Euler:}\markdownRendererParagraphSeparator
{}$$\markdownRendererSoftLineBreak
{}x^{\varphi(m)} \equiv 1 \pmod{m}, \quad x \text{ coprimo con } m\markdownRendererSoftLineBreak
{}$$\markdownRendererParagraphSeparator
{}Si $m$ es primo ($\varphi(m) = m - 1$):\markdownRendererInterblockSeparator
{}
\markdownRendererSectionEnd \markdownRendererSectionBegin
\markdownRendererHeadingThree{Pequeño Teorema de Fermat}\markdownRendererInterblockSeparator
{}$$\markdownRendererSoftLineBreak
{}x^{m - 1} \equiv 1 \pmod{m}\markdownRendererSoftLineBreak
{}$$\markdownRendererParagraphSeparator
{}
\markdownRendererSectionEnd 
\markdownRendererSectionEnd \markdownRendererSectionBegin
\markdownRendererHeadingTwo{Inversos Multiplicativos Modulares}\markdownRendererInterblockSeparator
{}El inverso multiplicativo de $x$ módulo $m$, $\mathrm{inv}_m(x)$, cumple:\markdownRendererParagraphSeparator
{}$$\markdownRendererSoftLineBreak
{}x \cdot \mathrm{inv}_m(x) \equiv 1 \pmod{m}\markdownRendererSoftLineBreak
{}$$\markdownRendererParagraphSeparator
{}Existe si y solo si $\gcd(x, m) = 1$, y se puede calcular con:\markdownRendererParagraphSeparator
{}$$\markdownRendererSoftLineBreak
{}\mathrm{inv}_m(x) = x^{\varphi(m) - 1} \pmod{m}\markdownRendererSoftLineBreak
{}$$\markdownRendererParagraphSeparator
{}Si $m$ es primo:\markdownRendererParagraphSeparator
{}$$\markdownRendererSoftLineBreak
{}\mathrm{inv}_m(x) = x^{m - 2} \pmod{m}\markdownRendererSoftLineBreak
{}$$\markdownRendererParagraphSeparator
{}\markdownRendererStrongEmphasis{Ejemplo:} $\mathrm{inv}_{17}(6) = 6^{15} \bmod 17 = 3$\markdownRendererParagraphSeparator
{}\markdownRendererInputFencedCode{./_markdown_main/2dc78622a1268e5a97bb5471766124f6.verbatim}{cpp}{cpp}\markdownRendererInterblockSeparator
{}$\clearpage$\markdownRendererInterblockSeparator
{}
\markdownRendererSectionEnd \markdownRendererSectionBegin
\markdownRendererHeadingTwo{Ecuaciones Diofánticas}\markdownRendererInterblockSeparator
{}Una ecuación diofántica tiene la forma:\markdownRendererParagraphSeparator
{}$$\markdownRendererSoftLineBreak
{}a x + b y = c, \quad a, b, c \in \mathbb{Z}\markdownRendererSoftLineBreak
{}$$\markdownRendererParagraphSeparator
{}Se puede resolver si $c$ es divisible por $\gcd(a,b)$.\markdownRendererHardLineBreak
{}Si $(x_0, y_0)$ es una solución, todas las soluciones son:\markdownRendererParagraphSeparator
{}$$\markdownRendererSoftLineBreak
{}\begin{cases}\markdownRendererSoftLineBreak
{}x = x_0 + k \dfrac{b}{\gcd(a,b)} \\\markdownRendererSoftLineBreak
{}y = y_0 - k \dfrac{a}{\gcd(a,b)} \\\markdownRendererSoftLineBreak
{}k \in \mathbb{Z}\markdownRendererSoftLineBreak
{}\end{cases}\markdownRendererSoftLineBreak
{}$$\markdownRendererParagraphSeparator
{}\markdownRendererStrongEmphasis{Ejemplo:} resolver $39x + 15y = 12$ $\gcd(39, 15) = 3 \mid 12$\markdownRendererParagraphSeparator
{}Usando el algoritmo de Euclides extendido:\markdownRendererParagraphSeparator
{}$39 \cdot 2 + 15 \cdot (-5) = 3 \implies 39 \cdot 8 + 15 \cdot (-20) = 12$\markdownRendererInterblockSeparator
{}\markdownRendererInputFencedCode{./_markdown_main/86307c4e753077e250ea378d789db39a.verbatim}{cpp}{cpp}\markdownRendererInterblockSeparator
{}$\clearpage$\markdownRendererInterblockSeparator
{}\markdownRendererSectionBegin
\markdownRendererHeadingThree{Teorema Chino del Resto}\markdownRendererInterblockSeparator
{}\markdownRendererStrongEmphasis{Sistema de congruencias:}\markdownRendererParagraphSeparator
{}$$\markdownRendererSoftLineBreak
{}\begin{cases}\markdownRendererSoftLineBreak
{}x \equiv a_1 \pmod{m_1} \\\markdownRendererSoftLineBreak
{}x \equiv a_2 \pmod{m_2} \\\markdownRendererSoftLineBreak
{}\vdots \\\markdownRendererSoftLineBreak
{}x \equiv a_n \pmod{m_n}\markdownRendererSoftLineBreak
{}\end{cases}\markdownRendererSoftLineBreak
{}$$\markdownRendererParagraphSeparator
{}con $m_i$ coprimos dos a dos.\markdownRendererParagraphSeparator
{}\markdownRendererStrongEmphasis{Solución general:}\markdownRendererParagraphSeparator
{}$$\markdownRendererSoftLineBreak
{}x = \sum_{k=1}^{n} a_k X_k , \mathrm{inv}_{m_k}(X_k)\markdownRendererSoftLineBreak
{}$$\markdownRendererParagraphSeparator
{}donde\markdownRendererParagraphSeparator
{}$$\markdownRendererSoftLineBreak
{}X_k = \frac{m_1 m_2 \cdots m_n}{m_k}\markdownRendererSoftLineBreak
{}$$\markdownRendererParagraphSeparator
{}\markdownRendererStrongEmphasis{Ejemplo:}\markdownRendererParagraphSeparator
{}$$\markdownRendererSoftLineBreak
{}\begin{cases}\markdownRendererSoftLineBreak
{}x \equiv 3 \pmod{5} \\\markdownRendererSoftLineBreak
{}x \equiv 4 \pmod{7} \\\markdownRendererSoftLineBreak
{}x \equiv 2 \pmod{3}\markdownRendererSoftLineBreak
{}\end{cases}\markdownRendererSoftLineBreak
{}\implies\markdownRendererSoftLineBreak
{}x = 263\markdownRendererSoftLineBreak
{}$$\markdownRendererParagraphSeparator
{}\markdownRendererStrongEmphasis{Infinitas soluciones:}\markdownRendererParagraphSeparator
{}$$\markdownRendererSoftLineBreak
{}x + m_1 m_2 \cdots m_n\markdownRendererSoftLineBreak
{}$$\markdownRendererInterblockSeparator
{}\markdownRendererInputFencedCode{./_markdown_main/ec69e9ebc473be347c5864afb079db77.verbatim}{cpp}{cpp}\markdownRendererInterblockSeparator
{}
\markdownRendererSectionEnd 
\markdownRendererSectionEnd \markdownRendererSectionBegin
\markdownRendererHeadingTwo{MEX}\markdownRendererInterblockSeparator
{}El \markdownRendererStrongEmphasis{MEX} ("minimum excludant") de un conjunto de números enteros es el \markdownRendererStrongEmphasis{mínimo número entero no presente en el conjunto}.\markdownRendererParagraphSeparator
{}Ejemplo: $\mathrm{mex}({0,1,3}) = 2$\markdownRendererInterblockSeparator
{}\markdownRendererInputFencedCode{./_markdown_main/35760a9625e11f9267eaf63a95669577.verbatim}{cpp}{cpp}\markdownRendererInterblockSeparator
{}
\markdownRendererSectionEnd 
\markdownRendererSectionEnd \markdownRendererSectionBegin
\markdownRendererHeadingOne{Juegos NIM}\markdownRendererInterblockSeparator
{}\markdownRendererUlBeginTight
\markdownRendererUlItem Hay (n) heaps con (x_i) palillos.\markdownRendererUlItemEnd 
\markdownRendererUlItem Jugadores alternan turnos.\markdownRendererUlItemEnd 
\markdownRendererUlItem En cada turno se remueven palillos de un heap no vacío.\markdownRendererUlItemEnd 
\markdownRendererUlItem Gana quien remueve el último palillo.\markdownRendererUlItemEnd 
\markdownRendererUlEndTight \markdownRendererInterblockSeparator
{}\markdownRendererSectionBegin
\markdownRendererHeadingTwo{Nim Sum}\markdownRendererInterblockSeparator
{}$$\markdownRendererSoftLineBreak
{}s = x_1 \oplus x_2 \oplus \dots \oplus x_n\markdownRendererSoftLineBreak
{}$$\markdownRendererInterblockSeparator
{}\markdownRendererUlBeginTight
\markdownRendererUlItem Si $(s = 0)$ estado perdedor.\markdownRendererUlItemEnd 
\markdownRendererUlItem Si $(s \neq 0)$ estado ganador.\markdownRendererUlItemEnd 
\markdownRendererUlEndTight \markdownRendererInterblockSeparator
{}
\markdownRendererSectionEnd \markdownRendererSectionBegin
\markdownRendererHeadingTwo{Movimiento óptimo}\markdownRendererInterblockSeparator
{}Para un estado ganador:\markdownRendererInterblockSeparator
{}\markdownRendererUlBeginTight
\markdownRendererUlItem Elegir heap $(i)$ tal que $(x_i \oplus s < x_i)$\markdownRendererUlItemEnd 
\markdownRendererUlItem Remover hasta que el heap tenga $(x_i \oplus s)$ palillos.\markdownRendererUlItemEnd 
\markdownRendererUlEndTight \markdownRendererParagraphSeparator
{}
\markdownRendererSectionEnd \markdownRendererSectionBegin
\markdownRendererHeadingTwo{Teorema de Sprague–Grundy}\markdownRendererInterblockSeparator
{}El Teorema de Sprague–Grundy generaliza la estrategia de Nim a cualquier juego que cumpla:\markdownRendererInterblockSeparator
{}\markdownRendererUlBeginTight
\markdownRendererUlItem Dos jugadores alternan turnos.\markdownRendererUlItemEnd 
\markdownRendererUlItem El juego consiste en estados, y los movimientos posibles dependen solo del estado.\markdownRendererUlItemEnd 
\markdownRendererUlItem El juego termina cuando un jugador no puede moverse.\markdownRendererUlItemEnd 
\markdownRendererUlItem El juego termina seguro eventualmente.\markdownRendererUlItemEnd 
\markdownRendererUlItem Los jugadores tienen información completa y no hay aleatoriedad.\markdownRendererUlItemEnd 
\markdownRendererUlEndTight \markdownRendererInterblockSeparator
{}$\clearpage$\markdownRendererInterblockSeparator
{}
\markdownRendererSectionEnd \markdownRendererSectionBegin
\markdownRendererHeadingTwo{Números de Grundy}\markdownRendererInterblockSeparator
{}Cada estado del juego se asocia con un \markdownRendererStrongEmphasis{número de Grundy} (G(s)), equivalente al tamaño de un heap de Nim.\markdownRendererParagraphSeparator
{}Si $(s)$ tiene movimientos hacia los estados $(s_1, s_2, \dots, s_n)$:\markdownRendererParagraphSeparator
{}$$\markdownRendererSoftLineBreak
{}\text{mex}({g_1,g_2,\markdownRendererEllipsis{},g_n})\markdownRendererSoftLineBreak
{}$$
\markdownRendererSectionEnd 
\markdownRendererSectionEnd \markdownRendererDocumentEnd